% Human-editable description for the MMS vortex verification test folder.
% Generated figures and tables are written by postProcess_mmsVortex.py under
% \CaseOutputDir and are included near the end of this section.

\paragraph{Purpose.}
This test verifies the generalized-vortex method-of-manufactured-solutions
(MMS) problem described by Kamojjala, Brannon, Sadeghirad, and Guilkey,
``Verification tests in solid mechanics,'' \emph{Engineering with Computers}
31:193--213, 2015, doi:10.1007/s00366-013-0342-x.  The test exercises
\texttt{SolidMechanics\_MPM} under a smooth large-deformation field containing
simple shear with superposed rotation, while keeping the analytical solution and
expected final state known.

\paragraph{Problem definition.}
The manufactured solution is posed on an annulus with inner radius $a=0.75$ and
outer radius $b=1.25$.  Material points move only in the circumferential
direction; the displacement and displacement gradients vanish at the inner and
outer radii, so the ring may also be embedded in a stationary square background
block.  The initial displacement and velocity are zero, the stop time is
$T=1$, and the temporal factor is zero again at $t=T$, so the exact final
configuration coincides with the initial configuration.  The material card is
the shared \texttt{pfw\_materials.hyperelasticMMS} parameterization of the
Neo-Hookean MMS model used for this test: density $1000$, Lam\'e parameter
$\lambda=577$, and shear modulus $\mu=384.6153846153846$ in the unit system of
the PFW input.

\paragraph{Numerical variants.}
The same human-readable PFW input is run twice from this folder.  One variant
uses CPDI particles with FLIP transfer, and one variant uses FMPM with
single-point B-spline particles.  The file \texttt{pfw\_input\_mmsVortex.py} is
the only PFW source input in this folder.  The folder-level \texttt{runTest}
script selects the numerical variant by setting documented template environment
variables before invoking \texttt{particleFileWriter.py}.  The template comments
list the concrete values used for the CPDI-FLIP and FMPM/singlePointBspline
runs, so users can copy the input and replace the defaults to make a standalone
single-case example.

\begin{figure}[h!]
  \centering
  \includegraphics[width=0.45\linewidth]{\CaseSourceDir/setup_vortex.png}\hfill
  \includegraphics[width=0.45\linewidth]{\CaseSourceDir/expected_deformation_kamojjala_fig11.png}
  \caption{Left: static setup diagram for the square computational domain,
  annular MMS motion region, and tracer rings.  Right: expected generalized
  vortex deformation pattern from Kamojjala et al.; hue is proportional to the
  initial angular coordinate.}
\end{figure}

\paragraph{Expected results.}
At $t=0$ the displacement field is zero.  Near $t=T/2$ the annulus should show a
smooth rotational vortex with maximum displacement inside the ring and zero
motion at the inner and outer radii.  At $t=T$ the exact displacement returns to
zero everywhere, so numerical residual displacement, shape distortion, or mesh
texture in the final configuration indicates integration, transfer, particle
support, or boundary-treatment error.  Kamojjala et al. used the generalized
vortex to reveal that older uGIMP-style MPM could become unstable in large
deformation while CPDI remained stable; this suite compares the current CPDI-FLIP
and FMPM/singlePointBspline implementations against the same manufactured
motion.

\paragraph{Metric.}
The quantitative metric is the average norm of the displacement of a fixed
array of Lagrangian tracer particles seeded in the annulus.  The post-processor
compares this history against the analytical displacement norm from the
manufactured vortex and reports peak and final RMS errors for both variants.
The suite also records Slurm timing metadata when available, and the input
requests current GEOS solver statistics with \texttt{writeStatistics="all"} so
run directories retain detailed solver-side diagnostics.

\paragraph{Generated outputs.}
The post-processor writes combined CSV data, JSON summaries, comparison plots,
and the VisIt particle/background images at $t=0$, $t=t_{stop}/2$, and
$t=t_{stop}$ when VisIt is available.

\IfFileExists{\CaseOutputDir/mms_vortex_results.tex}{%
  \input{\CaseOutputDir/mms_vortex_results.tex}%
}{%
  \begin{center}\fbox{MMS vortex post-processing products have not been generated.}\end{center}%
}
